% ============================================================
%  Поглавље 3: Проблем MaxCut
% ============================================================
\section{Проблем MaxCut}
\label{sec:maxcut}

% --- Формална дефиниција ---
% Дефинишите граф G=(V,E) са тежинама. Дефинишите рез (бисекцију)
% скупа V у два скупа S и V\S. Циљ: максимизовати тежину
% ивица које прелазе рез.

\subsection{Формална дефиниција}

\begin{definition}[MaxCut]
Нека је $G = (V, E, w)$ неусмерени граф са тежинама $w : E \to \mathbb{R}_{+}$.
\emph{Рез} је партиција $V = S \cup \bar{S}$, $S \cap \bar{S} = \emptyset$.
Проблем \textsc{MaxCut} тражи партицију која максимизује
\[
  \text{cut}(S) = \sum_{\{u,v\} \in E,\; u \in S,\; v \in \bar{S}} w(u,v).
\]
\end{definition}

% --- Пример на малом графу ---
% Нацртајте граф са 5--6 чворова (TikZ), прикажите оптималан рез
% и вредност циља. Ово ће послужити и у каснијим поглављима
% (Z3 кодирање, QAOA).

\subsection{Илустративни пример}

% TODO: Нацртати TikZ граф, означити оптималан рез испрекиданом линијом.
% Препоручени пример: K4 граф (потпуни граф на 4 чвора), MaxCut = 4.

\subsection{НП-тешкоћа MaxCut-а}

% --- Редукција из 3-SAT (или NAE-3SAT) ---
% Скицирајте доказ НП-тешкоћа MaxCut-а полиномском редукцијом
% из НП-потпуног проблема (нпр. NAE-3SAT → MaxCut).
% Укључите референцу на Карп 1972.

\begin{theorem}[Карп, 1972~\cite{karp1972}]
Проблем \textsc{MaxCut} је \textsf{NP}-тешки.
\end{theorem}

% TODO: Скица доказа редукцијом.

\subsection{Класични апроксимациони алгоритми}

% --- Гоеманс-Вилијамсон алгоритам ---
% Опишите SDP релаксацију и хиперраванску рандомизацију.
% Наведите гаранцију апроксимације: ≥ 0.878 · OPT.
% Ово мотивише питање: може ли квантни приступ да уради боље?

Гоеманс и Вилијамсон~\cite{goemans1995} предложили су алгоритам заснован
на семидефинитном програмирању (SDP) са гаранцијом апроксимације:
\[
  \text{cut}(S_{\text{alg}}) \;\geq\; 0.878 \cdot \text{OPT}.
\]

% TODO: Опис SDP релаксације и хиперраванске рандомизације.
% Напомена: под условом P ≠ NP и UGC хипотезе, ово је оптимална
% апроксимација за полиномске алгоритме.
